%%===================================================================
%% Lecture 21 -- The Standard Model Completed:
%%               Spectrum, Neutral/Charged/Higgs Currents,
%%               Accidental Symmetries, Evidence for Color,
%%               and Beyond the Standard Model
%% 77609 -- Introduction to Elementary Particles
%% Date: 25/06/2026  (final lecture)
%% Lecturer: Prof. Yonit Hochberg
%%===================================================================

\section{Lecture 21 -- The Standard Model Completed: Currents, Color, Accidental Symmetries, and Beyond}\label{sec:lec21}
\begin{center}
\begin{tabular}{ll}
  \textbf{Date:}\quad 25/06/2026 \quad(\emph{final lecture}) & \\
\end{tabular}
\end{center}
\hrule\vspace{1.5em}

%%------------------------------------------------------------------
\subsection{Overview}\label{subsec:lec21:overview}

This is the closing lecture of the course. In
Lectures~17--20 we built the full Standard Model gauge group
$G_{\SM}=SU(3)_C\times SU(2)_L\times U(1)_Y$, broke it spontaneously to
$SU(3)_C\times U(1)_{\rm EM}$, and discovered the CKM matrix in the
quark Yukawa sector. Today we \emph{harvest the physics} of that
construction: we read off every interaction the SM predicts, classify
them, confront them with experiment, and finally step past the model to
the open questions that define modern high-energy research.

The logical chain is:
\begin{enumerate}
  \item \textbf{Fermion spectrum} --- masses from the Higgs vev, the full
        mass table, and the (theoretically) massless neutrinos.
  \item \textbf{Neutral-current ($Z$) interactions} --- the master
        coupling $\propto (T^3-\sin^2\theta_W\,Q)$; universal, diagonal,
        chiral.
  \item \textbf{Charged-current ($W$) interactions} --- trivial for
        leptons, but for quarks the CKM matrix makes them
        \emph{non-universal} and \emph{non-diagonal}; only left-handed
        fields couple.
  \item \textbf{Higgs interactions} --- couplings $\propto$ mass; why the
        Higgs decays dominantly to $b\bar b$; why there is no tree-level
        $hgg$ or $h\gamma\gamma$.
  \item \textbf{Accidental symmetries} --- baryon and lepton number,
        and why the proton is stable.
  \item \textbf{CKM magnitudes} --- the hierarchical Wolfenstein
        structure.
  \item \textbf{Evidence for color} --- the $R$-ratio of
        $e^+e^-$ annihilation directly counts $N_c=3$.
  \item \textbf{Weak-process taxonomy} --- leptonic, semileptonic,
        non-leptonic.
  \item \textbf{Beyond the Standard Model} --- naturalness, dark matter,
        baryogenesis, the flavor puzzle.
\end{enumerate}

%%==================================================================
\subsection{The Fermion Spectrum}\label{subsec:lec21:spectrum}

\subsubsection*{Physical Motivation}

In Lecture~20 we counted the \emph{boson} spectrum (massive $W^\pm,Z$;
massless $\gamma$ and $8$ gluons). We now complete the picture with the
\emph{fermion} masses. They are not free constants written into the
Lagrangian --- gauge invariance forbade bare masses --- so every one of
them is generated when the neutral Higgs component acquires its vacuum
expectation value $\langle\phi^0\rangle = v/\sqrt2$.

\subsubsection*{Masses from the Yukawa couplings}

Inserting $\langle\phi^0\rangle=v/\sqrt2$ into the Yukawa terms of
Lecture~20, each diagonalized Yukawa eigenvalue becomes a mass:
\begin{equation}
  \boxed{
    m_{\ell} = y_\ell\,\frac{v}{\sqrt2},\quad
    m_{u_i} = y_{u_i}\,\frac{v}{\sqrt2},\quad
    m_{d_i} = y_{d_i}\,\frac{v}{\sqrt2},
  }
  \label{eq:lec21:fermion-masses}
\end{equation}
where $\ell\in\{e,\mu,\tau\}$, $u_i\in\{u,c,t\}$, $d_i\in\{d,s,b\}$, and
$y$ are the (real, diagonal) Yukawa eigenvalues. Because the SM contains
\emph{no} right-handed neutrino and (at the renormalizable level) no way
to write a neutrino mass, the three neutrinos are predicted massless:
\begin{equation}
  \boxed{m_{\nu_e}=m_{\nu_\mu}=m_{\nu_\tau}=0
        \quad(\text{in the SM}).}
  \label{eq:lec21:massless-nu}
\end{equation}

\textit{Significance:} A single number, $v\approx 246~\text{GeV}$, sets
the overall scale of \emph{every} fermion mass; the enormous spread of
masses is entirely the spread of the dimensionless Yukawa couplings $y$.

\nt{Pitfall: Eq.~\eqref{eq:lec21:massless-nu} is the SM
\emph{prediction}, and it is \emph{wrong} --- neutrino oscillations show
$m_\nu\neq0$. This is the one place where the renormalizable SM
demonstrably fails, and it is the first item on the ``beyond the SM''
list in Sec.~\ref{subsec:lec21:bsm}.}

\paragraph{The measured spectrum.}
Collecting the experimental values (used repeatedly as a reference for
``what can decay to what''):
\begin{center}
\renewcommand{\arraystretch}{1.25}
\begin{tabular}{@{}llll@{}}
  \toprule
  Generation & Charged lepton & Up-type quark & Down-type quark\\
  \midrule
  1 & $m_e\approx 0.511~\text{MeV}$    & $m_u\approx 1\text{--}3~\text{MeV}$  & $m_d\approx 4\text{--}6~\text{MeV}$\\
  2 & $m_\mu\approx 105.66~\text{MeV}$ & $m_c\approx 1.3~\text{GeV}$          & $m_s\approx 100~\text{MeV}$\\
  3 & $m_\tau\approx 1776.82~\text{MeV}$ & $m_t\approx 172.9~\text{GeV}$      & $m_b\approx 4.9~\text{GeV}$\\
  \bottomrule
\end{tabular}
\end{center}
The top quark is the heaviest fundamental particle in the SM, and the
Higgs sits at $m_h\approx 125~\text{GeV}$. These two facts together
($2m_t>m_h$, but $2m_b<m_h$) decide the dominant Higgs decay below.

%%==================================================================
\subsection{Categories of Interaction}\label{subsec:lec21:categories}

Once the spectrum is fixed, every interaction is \emph{read off} from
the covariant derivative and the Yukawa terms. They fall into groups:
\begin{itemize}
  \item \textbf{Already derived} (nothing new): QED of leptons, QED of
        quarks, and QCD of quarks --- identical in form to what we built
        for the Lepton Standard Model and for QCD.
  \item \textbf{To work out now}: the couplings to the massive bosons
        $Z$, $W^\pm$, and $h$, which carry the genuinely electroweak
        physics.
\end{itemize}
We treat the three new cases in turn.

%%==================================================================
\subsection{Neutral-Current (\texorpdfstring{$Z$}{Z}) Interactions}\label{subsec:lec21:nc}

\subsubsection*{Conceptual Background}

The ``neutral current'' is the interaction mediated by the $Z$ boson.
After electroweak symmetry breaking the $Z$ couples to every fermion
through a single universal structure, weighted by the fermion's weak
isospin $T^3$ and electric charge $Q$:
\begin{equation}
  \boxed{
    \lagr_{\rm NC}
    = \frac{e}{\sin\theta_W\cos\theta_W}
      \sum_i \bar\psi_i
      \left(T^3_i - \sin^2\theta_W\,Q_i\right)
      \slashed{Z}\,\psi_i ,
  }
  \label{eq:lec21:nc}
\end{equation}
where $\theta_W$ is the weak mixing angle, $e=g\sin\theta_W$ is the
electric charge, the sum runs over \emph{all} fermions $\psi_i$, and
$\slashed Z=\gamma^\mu Z_\mu$. The prefactor can equivalently be written
$g/\cos\theta_W$.

\textit{Significance:} The presence of $T^3$ inside
Eq.~\eqref{eq:lec21:nc} means the $Z$ \emph{cares about chirality}:
left-handed fields (doublets, $T^3=\pm\tfrac12$) and right-handed fields
(singlets, $T^3=0$) couple differently. The $Z$ is a \emph{chiral}
probe.

\subsubsection*{Reading off a coefficient}

For each fermion we only need two numbers: its $T^3$ (the upper
component of a doublet has $+\tfrac12$, the lower has $-\tfrac12$, a
singlet has $0$) and its electric charge $Q=T^3+Y$. The coefficient is
then $T^3-\sin^2\theta_W\,Q$. Recall the representations from
Lecture~20:
\[
  Q_L=(\mathbf3,\mathbf2)_{+1/6},\quad
  u_R=(\mathbf3,\mathbf1)_{+2/3},\quad
  d_R=(\mathbf3,\mathbf1)_{-1/3},\quad
  L_L=(\mathbf1,\mathbf2)_{-1/2},\quad
  e_R=(\mathbf1,\mathbf1)_{-1}.
\]

\ex[ex:lec21:zcoeff]{Worked Example: $Z$-Coupling Coefficients}{
Compute the coefficient $c_\psi = T^3_\psi - \sin^2\theta_W\,Q_\psi$ for
four representative fields. Abbreviate $s_W^2\equiv\sin^2\theta_W$.
\begin{itemize}
  \item \textbf{Left-handed charged lepton $\ell_L$} (lower component of
        $L_L$): $T^3=-\tfrac12$, $Q=-1$, so
        \[
          c_{\ell_L} = -\tfrac12 - s_W^2(-1) = -\tfrac12 + s_W^2 .
        \]
  \item \textbf{Right-handed charged lepton $e_R$} (singlet):
        $T^3=0$, $Q=-1$, so
        \[
          c_{e_R} = 0 - s_W^2(-1) = +\,s_W^2 .
        \]
  \item \textbf{Left-handed neutrino $\nu_L$} (upper component of $L_L$):
        $T^3=+\tfrac12$, $Q=0$, so
        \[
          c_{\nu_L} = +\tfrac12 - s_W^2(0) = +\tfrac12 .
        \]
  \item \textbf{Left-handed down quark $d_L$} (lower component of $Q_L$):
        $T^3=-\tfrac12$, $Q=-\tfrac13$, so
        \[
          c_{d_L} = -\tfrac12 - s_W^2\!\left(-\tfrac13\right)
                  = -\tfrac12 + \tfrac13 s_W^2 .
        \]
\end{itemize}
Each coupling in the Lagrangian is then
$\dfrac{e}{s_W c_W}\,c_\psi\,\bar\psi\,\slashed Z\,\psi$ with
$c_W=\cos\theta_W$. Note $c_{\ell_L}\neq c_{e_R}$: the left- and
right-handed electron couple to the $Z$ with \emph{different} strengths
--- the hallmark of a chiral coupling.
}

\subsubsection*{Three defining properties}

\clm[clm:lec21:nc-props]{The Neutral Current is Universal, Diagonal, and Chiral}{
\begin{itemize}
  \item \textbf{Universal:} within a given sector (e.g.\ all left-handed
        charged leptons, or all left-handed up-type quarks) the coupling
        is \emph{generation-independent} --- $e_L,\mu_L,\tau_L$ all
        couple identically.
  \item \textbf{Diagonal:} the $Z$ never connects two \emph{different}
        flavors; there is no $\bar e\,\slashed Z\,\mu$ vertex
        (no tree-level flavor-changing neutral current).
  \item \textbf{Chiral:} the coupling depends on $T^3$, hence on whether
        the field is left- or right-handed.
\end{itemize}
All three properties are \emph{confirmed by experiment}.
}

\textit{Significance:} ``Universal $+$ diagonal $+$ chiral'' is a sharp,
falsifiable prediction of the electroweak gauge structure; the absence
of tree-level flavor-changing neutral currents (FCNCs) is one of its
most successful consequences.

\subsubsection*{Method: Writing any $Z$ coupling}
\begin{enumerate}
  \item Identify the fermion's $SU(2)_L$ representation; read $T^3$
        ($+\tfrac12$ upper doublet, $-\tfrac12$ lower doublet, $0$
        singlet).
  \item Compute $Q=T^3+Y$.
  \item Form $c_\psi=T^3-\sin^2\theta_W\,Q$.
  \item The interaction term is
        $\dfrac{e}{\sin\theta_W\cos\theta_W}\,c_\psi\,
        \bar\psi\,\slashed Z\,\psi$.
\end{enumerate}

\nt{Pitfall: $T^3$ is the eigenvalue of $\sigma^3/2$, \emph{not} of
$\sigma^3$. The upper/lower components of a doublet carry
$T^3=\pm\tfrac12$ (not $\pm1$). Forgetting the factor $\tfrac12$ is the
single most common error in computing $Z$ couplings.}

%%==================================================================
\subsection{Charged-Current (\texorpdfstring{$W^\pm$}{W}) Interactions}\label{subsec:lec21:cc}

\subsubsection*{Conceptual Background}

The ``charged current'' is mediated by the $W^\pm$, which carries
electric charge $\pm1$. Charge conservation at the vertex forces the
$W$ to connect two \emph{different} fermions whose charges differ by one
unit: a charged lepton and its neutrino, or an up-type and a down-type
quark. Because the right-handed fields are $SU(2)_L$ singlets, only
\emph{left-handed} fields couple to the $W$.

\subsubsection*{Leptons: the easy case}

For leptons the mass basis and the interaction (weak) basis coincide
--- there is no leptonic mixing because the neutrinos are massless and
only $Y^e$ needs diagonalizing. The coupling is therefore
flavor-diagonal and generation-universal:
\begin{equation}
  \boxed{
    \lagr_{\rm CC}^{\rm lept}
    = -\frac{g}{\sqrt2}
      \left(
        \bar\nu_{eL}\,\slashed{W}^+ e_L
        + \bar\nu_{\mu L}\,\slashed{W}^+ \mu_L
        + \bar\nu_{\tau L}\,\slashed{W}^+ \tau_L
      \right) + \text{h.c.}
  }
  \label{eq:lec21:cc-lept}
\end{equation}
Each $W$ links a charged lepton only to its \emph{own} neutrino: there
is no $\bar\nu_\mu\,\slashed W\,e$ term. This is lepton flavor
universality.

\textit{Significance:} The $W$ vertex is what makes one neutrino
``belong'' to one charged lepton --- it operationally \emph{defines}
electron-, muon-, and tau-number.

\subsubsection*{Quarks: the CKM matrix enters}

For quarks the interaction basis is \emph{not} the mass basis. Working
in the down-quark mass basis, the $W$ couples the up-type mass
eigenstates $(u_L,c_L,t_L)$ to CKM-rotated down-type fields. Writing the
mass eigenstates as $(u_L,c_L,t_L)^T = V_{\rm CKM}\,(u^d_L,c^d_L,t^d_L)^T$
where the superscript $d$ denotes the partner sitting above $(d,s,b)$ in
the down basis, the charged-current Lagrangian becomes
\begin{equation}
  \boxed{
    \lagr_{\rm CC}^{\rm quark}
    = -\frac{g}{\sqrt2}
      \begin{pmatrix}\bar u_L & \bar c_L & \bar t_L\end{pmatrix}
      V_{\rm CKM}\,\slashed{W}^+
      \begin{pmatrix} d_L \\ s_L \\ b_L \end{pmatrix}
    + \text{h.c.}
  }
  \label{eq:lec21:cc-quark}
\end{equation}

\clm[clm:lec21:cc-props]{The Quark Charged Current is Left-Handed, Non-Universal, and Non-Diagonal}{
\begin{itemize}
  \item \textbf{Only left-handed fields couple} --- the right-handed
        quarks are $SU(2)_L$ singlets, so they never see the $W$.
  \item \textbf{Non-universal and non-diagonal} --- $V_{\rm CKM}$ is not
        proportional to the identity, so the $W$ connects different
        generations (e.g.\ an up quark to a strange quark) with
        amplitude $(V_{\rm CKM})_{us}$.
\end{itemize}
}

\textit{Significance:} The CKM matrix lives \emph{entirely} in the quark
charged current. It is the sole source of quark flavor change (and of CP
violation) in the SM.

\subsubsection*{Decay rates and the color factor}

When a $W$ decays to a quark pair $u_i\bar d_j$, the amplitude is
proportional to the relevant CKM element, so the rate goes as its
modulus squared. There is also a multiplicity factor of $3$ because each
quark is an $SU(3)_C$ triplet (three colors in the final state):
\begin{equation}
  \boxed{
    \Gamma(W^+\to u_i\,\bar d_j)\ \propto\ 3\,\bigl|(V_{\rm CKM})_{ij}\bigr|^2 .
  }
  \label{eq:lec21:w-decay}
\end{equation}
By contrast, the leptonic rate $\Gamma(W^+\to \ell^+\nu_\ell)\propto 1$
(no CKM, no color factor). Kinematically the $W$ \emph{cannot} decay to
a final state containing the top quark, because $m_t>m_W$ --- the top is
too heavy to be produced in a $W$ decay at rest.

\ex[ex:lec21:wdecay]{Worked Example: Comparing $W$ Decay Channels}{
A leptonic channel such as $W^+\to e^+\nu_e$ has amplitude
$\propto 1$ (dropping the common $g/\sqrt2$). For the hadronic channel
$W^+\to u\bar d$, the amplitude carries the CKM element $V_{ud}$ and the
rate carries an extra color factor $3$:
\[
  \frac{\Gamma(W^+\to u\bar d)}{\Gamma(W^+\to e^+\nu_e)}
  \approx 3\,|V_{ud}|^2 \approx 3(0.974)^2 \approx 2.85 .
\]
Summing over all kinematically allowed quark pairs (and using CKM
unitarity, $\sum_j|V_{ij}|^2=1$) gives roughly $3$ for each accessible
up-type row --- the origin of the well-known $\sim$ $1:1:1:3:3$ ratio of
$e\nu:\mu\nu:\tau\nu:(\text{hadrons from }u):(\text{hadrons from }c)$
branching fractions. The factor $3$ is a direct count of color.
}

\nt{Pitfall: the top does not appear in $W$ decays not because its CKM
element is small, but because it is \emph{kinematically forbidden}
($m_t>m_W$). Suppression by a small $|V_{ij}|$ and exclusion by
kinematics are different effects --- do not conflate them.}

%%==================================================================
\subsection{Higgs-Boson Interactions}\label{subsec:lec21:higgs}

\subsubsection*{Conceptual Background}

The same Yukawa coupling that gives a fermion its mass also controls its
coupling to the physical Higgs $h$. Hence the Higgs coupling to a
fermion is \emph{proportional to that fermion's mass} and is
\emph{flavor-diagonal}:
\begin{equation}
  \boxed{
    \lagr_{hf\bar f}
    = -\sum_f \frac{m_f}{v}\,h\,\bar f_L f_R + \text{h.c.},
    \qquad f\in\{e,\mu,\tau,u,d,s,c,b,t\}.
  }
  \label{eq:lec21:hff}
\end{equation}
This follows because, after symmetry breaking,
$\phi^0 \to (v+h)/\sqrt2$, so the Yukawa term $y_f\bar f_L\phi^0 f_R$
produces both the mass $m_f=y_f v/\sqrt2$ and the coupling
$(y_f/\sqrt2)\,h\,\bar f_L f_R = (m_f/v)\,h\,\bar f_L f_R$.

\textit{Significance:} ``Coupling $\propto$ mass'' is the signature
prediction of the Higgs mechanism; measuring Higgs couplings to
different fermions and finding they scale with mass is the central test
of the Higgs sector.

\subsubsection*{The full Higgs Lagrangian}

Collecting the self-interactions and the couplings to the massive gauge
bosons (each $\propto m_V^2$, because gauge-boson masses also come from
the vev):
\begin{equation}
  \begin{aligned}
    \lagr_h
    ={}& \tfrac12(\partial_\mu h)^2 - \tfrac12 m_h^2 h^2
       - \frac{m_h^2}{2v}\,h^3 - \frac{m_h^2}{8v^2}\,h^4 \\
     &+ m_W^2\,W^+_\mu W^{-\mu}
        \left(\frac{2h}{v}+\frac{h^2}{v^2}\right)
      + \tfrac12 m_Z^2\,Z_\mu Z^{\mu}
        \left(\frac{2h}{v}+\frac{h^2}{v^2}\right)
      - \sum_f \frac{m_f}{v}\,h\,\bar f_L f_R + \text{h.c.}
  \end{aligned}
  \label{eq:lec21:higgs-lagr}
\end{equation}

\subsubsection*{Which decay dominates?}

The Higgs couples \emph{most strongly} to the heaviest fermion, the top.
But a \emph{decay} $h\to f\bar f$ requires $2m_f<m_h$. Since
$2m_t\approx 346~\text{GeV} > m_h\approx 125~\text{GeV}$, the top channel
is closed. Among the pair-producible fermions, the heaviest is the
bottom quark, so:
\begin{equation}
  \boxed{
    h\ \text{decays dominantly to}\ b\bar b
    \quad(\text{heaviest pair that satisfies } 2m_f<m_h).
  }
  \label{eq:lec21:hbb}
\end{equation}

\textit{Significance:} The dominant Higgs decay is a competition between
``stronger coupling to heavier fermions'' and ``kinematic access'';
$b\bar b$ wins because the top is just out of reach.

\subsubsection*{No tree-level $hgg$ or $h\gamma\gamma$}

There is \emph{no} tree-level coupling of the Higgs to gluons or
photons. The reason is symmetric on both sides:
\begin{itemize}
  \item Gluons and the photon are massless (gauge bosons of the
        \emph{unbroken} $SU(3)_C$ and $U(1)_{\rm EM}$). The Higgs couples
        to gauge bosons \emph{through their mass}, so a massless boson
        gets no tree-level $h$ coupling.
  \item Equivalently, $h$ is color-neutral and electrically neutral, so
        it cannot couple at tree level to the carriers of color or
        charge.
\end{itemize}
These processes \emph{do} occur at one loop. For example
$h\to\gamma\gamma$ and $h\to gg$ proceed through a fermion
\emph{triangle diagram}, dominated by the top quark (strongest coupling)
running in the loop --- even though a real top cannot be produced, a
\emph{virtual} top can. The $h\to\gamma\gamma$ loop is the channel
through which the Higgs was originally discovered.

\begin{figure}[ht]
\centering
\begin{tikzpicture}[>=Stealth, line join=round]
  % incoming Higgs (dashed)
  \draw[dashed, thick] (-2.6,0) -- (-1.2,0) node[midway, above]{$h$};
  % the fermion triangle
  \coordinate (A) at (-1.2,0);
  \coordinate (B) at (0.6,1.0);
  \coordinate (C) at (0.6,-1.0);
  \draw[->, thick] (A) -- (B);
  \draw[->, thick] (B) -- (C);
  \draw[->, thick] (C) -- (A);
  \node at (-0.55,0) {$t$};
  % outgoing photons (wavy)
  \draw[thick, decorate, decoration={snake, amplitude=1.1pt, segment length=6pt}]
       (B) -- (2.4,1.4) node[right]{$\gamma$};
  \draw[thick, decorate, decoration={snake, amplitude=1.1pt, segment length=6pt}]
       (C) -- (2.4,-1.4) node[right]{$\gamma$};
\end{tikzpicture}
\caption{The top-quark triangle loop generating $h\to\gamma\gamma$
(and, with gluons on the right, $h\to gg$). No tree-level vertex
exists; the coupling is induced by a virtual heavy fermion.}
\label{fig:lec21:triangle}
\end{figure}

\nt{Pitfall: ``the Higgs couples most strongly to the top'' and ``the
Higgs cannot decay to the top'' are both true and not in conflict. The
top dominates \emph{loop-induced} processes (as a virtual particle) but
is forbidden as a \emph{real} decay product by $2m_t>m_h$.}

%%==================================================================
\subsection{Accidental Symmetries and Proton Stability}\label{subsec:lec21:accidental}

\subsubsection*{Conceptual Background}

An \emph{accidental symmetry} is a global symmetry of the Lagrangian
that we did not impose --- it appears automatically once we write down
the most general renormalizable terms consistent with the gauge
symmetry and field content. We diagnose them by the standard trick:
ask what symmetry survives when the Yukawa couplings are switched on and
off.

\subsubsection*{Yukawas off: a large flavor symmetry}

With all Yukawas set to zero, only the kinetic terms remain. Each of the
five fermion representations comes in three identical generations, so we
may rotate each representation independently in generation space by a
$U(3)$:
\begin{equation}
  \boxed{
    G_{\rm flavor}^{(Y=0)} = U(3)^5
    \;=\; U(3)_Q\times U(3)_u\times U(3)_d\times U(3)_L\times U(3)_e .
  }
  \label{eq:lec21:u3-5}
\end{equation}

\textit{Significance:} The covariant derivative is blind to generation,
so before the Yukawas the SM has an enormous global flavor symmetry. The
Yukawa couplings are the \emph{only} terms that break it.

\subsubsection*{Yukawas on: what survives}

Turning the Yukawas back on breaks $U(3)^5$ down dramatically:
\begin{itemize}
  \item \textbf{Leptons.} Each charged-lepton mass is distinct, so only
        independent \emph{phase} rotations survive --- one per
        generation. These are the individual lepton numbers
        $U(1)_e\times U(1)_\mu\times U(1)_\tau$.
  \item \textbf{Quarks.} Here something subtler happens. The
        \emph{same} doublet $Q_L$ appears in \emph{both} the up Yukawa
        ($Y^u$) and the down Yukawa ($Y^d$). One cannot rotate the
        up-type and down-type quarks by independent phases without
        spoiling one of the two Yukawa terms. The only surviving phase
        rotation is the \emph{common} one acting on all quarks
        simultaneously --- baryon number.
\end{itemize}
The surviving accidental symmetry is therefore
\begin{equation}
  \boxed{
    G_{\rm flavor}^{\rm SM}
    = U(1)_B\times U(1)_e\times U(1)_\mu\times U(1)_\tau,
  }
  \label{eq:lec21:accidental-sm}
\end{equation}
where under baryon number $U(1)_B$ every quark carries charge
$+\tfrac13$ (so a three-quark baryon has $B=1$) and all other fields are
neutral.

\textit{Significance:} Baryon number and the three lepton numbers are
\emph{not} imposed --- they emerge automatically. They are exactly
conserved by the renormalizable SM.

\subsubsection*{Why the proton is stable}

\thm[thm:lec21:proton]{Stability of the Lightest Charged State}{
The lightest particle carrying a conserved (accidental) charge is
\emph{stable}: it has no lighter state of the same charge to decay into.
The proton, being the lightest state with $B=1$, is therefore stable:
\begin{equation}
  \boxed{\text{proton stable} \iff U(1)_B \text{ conserved}.}
  \label{eq:lec21:proton}
\end{equation}
}

The proton ($uud$) is the lightest baryon-number-carrying state, so
$U(1)_B$ conservation forbids proton decay --- consistent with the fact
that proton decay has never been observed. However, $U(1)_B$ is only an
\emph{accidental} symmetry of the \emph{renormalizable} Lagrangian:
higher-dimensional (non-renormalizable) operators, expected from
physics beyond the SM, can violate $B$ and \emph{would} make the proton
decay. Observing proton decay would thus be direct evidence for new
physics.

\textit{Significance:} The longevity of the proton is not a brute fact
but a consequence of an accidental symmetry --- and its eventual decay,
if seen, would be a window onto a deeper Lagrangian.

%%==================================================================
\subsection{The CKM Matrix: Measured Magnitudes}\label{subsec:lec21:ckm-values}

\subsubsection*{Physical Motivation}

In Lecture~20 we defined $V_{\rm CKM}=V_{uL}V_{dL}^\dagger$ abstractly.
Here are its measured magnitudes, which reveal a striking hierarchical
pattern.

The measured magnitudes are approximately
\begin{equation}
  |V_{\rm CKM}| \approx
  \begin{pmatrix}
    |V_{ud}| & |V_{us}| & |V_{ub}| \\
    |V_{cd}| & |V_{cs}| & |V_{cb}| \\
    |V_{td}| & |V_{ts}| & |V_{tb}|
  \end{pmatrix}
  \approx
  \begin{pmatrix}
    0.974 & 0.225 & 0.0035 \\
    0.225 & 0.973 & 0.041 \\
    0.0087 & 0.041 & 0.999
  \end{pmatrix}.
  \label{eq:lec21:ckm-numbers}
\end{equation}
The matrix is nearly the identity, with small off-diagonal elements that
shrink rapidly away from the diagonal. Introducing the Cabibbo parameter
$\lambda\equiv|V_{us}|\approx 0.225$, the structure is captured by the
Wolfenstein parametrization:
\begin{equation}
  \boxed{
    V_{\rm CKM} \approx
    \begin{pmatrix}
      1 & \lambda & \lambda^3 \\
      \lambda & 1 & \lambda^2 \\
      \lambda^3 & \lambda^2 & 1
    \end{pmatrix},
    \qquad \lambda\approx 0.225 .
  }
  \label{eq:lec21:wolfenstein}
\end{equation}

\textit{Significance:} Same-generation transitions are unsuppressed,
$1\leftrightarrow2$ transitions cost $\lambda$, $2\leftrightarrow3$ cost
$\lambda^2$, and $1\leftrightarrow3$ cost $\lambda^3$. This hierarchy
directly controls which $W$-mediated quark decays are fast and which are
rare (rates $\propto|V_{ij}|^2$).

\nt{Parameter counting: a $3\times3$ unitary matrix has $9$ real
parameters; $5$ relative quark phases can be removed by field
redefinitions, leaving $4$ physical parameters --- \emph{three mixing
angles and one CP-violating phase}. That irremovable phase, impossible
with only two generations, is why Kobayashi and Maskawa predicted a
third generation, and why this matrix earned the 2008 Nobel Prize.}

%%==================================================================
\subsection{Evidence for Color: the \texorpdfstring{$R$}{R}-Ratio}\label{subsec:lec21:color}

\subsubsection*{Physical Motivation}

How do we \emph{know} there are exactly $N_c=3$ colors? One of the
cleanest pieces of evidence comes from comparing two cross sections in
$e^+e^-$ annihilation: production of muons versus production of hadrons
(quarks). The ratio counts colors directly.

\subsubsection*{Which diagram dominates}

At tree level, $e^+e^-\to f\bar f$ can proceed through an $s$-channel
photon, $Z$, or Higgs:
\begin{itemize}
  \item \textbf{Photon:} propagator $\sim 1/q^2$ (massless photon).
  \item \textbf{$Z$:} propagator $\sim 1/(q^2-m_Z^2)$ --- suppressed by
        the large $m_Z$ at energies $q^2\ll m_Z^2$.
  \item \textbf{Higgs:} coupling $\propto m_e$ (the electron Yukawa) ---
        suppressed by the tiny electron mass.
\end{itemize}
At energies well below $m_Z$, the \emph{photon exchange dominates}
completely.

\subsubsection*{The ratio counts color}

Because the initial state ($e^+e^-$) and the muon have the same electric
charge ($|Q|=1$), those factors cancel in the ratio, leaving only the
sum over the quarks' squared charges, times the number of colors:
\begin{equation}
  \boxed{
    R \equiv
    \frac{\sigma(e^+e^-\to \text{hadrons})}{\sigma(e^+e^-\to\mu^+\mu^-)}
    = N_c\sum_q Q_q^2 ,
  }
  \label{eq:lec21:Rratio}
\end{equation}
where the sum runs over all quark flavors light enough to be
pair-produced at the given energy.

\ex[ex:lec21:Rratio]{Worked Example: $R$ Below the Top Threshold}{
At energies below the top mass ($\sqrt s < 2m_t$) but above the $b\bar b$
threshold, the accessible quarks are $u,d,s,c,b$:
\begin{itemize}
  \item up-type ($Q=+\tfrac23$): $u,c$ $\Rightarrow$
        $2\times\left(\tfrac23\right)^2 = \tfrac89$;
  \item down-type ($Q=-\tfrac13$): $d,s,b$ $\Rightarrow$
        $3\times\left(\tfrac13\right)^2 = \tfrac39 = \tfrac13$.
\end{itemize}
So $\sum_q Q_q^2 = \tfrac89+\tfrac39 = \tfrac{11}{9}$. Then:
\[
  R_{\text{no color}} = \frac{11}{9}\approx 1.22,
  \qquad
  R_{\text{with } N_c=3} = 3\times\frac{11}{9} = \frac{33}{9}\approx 3.67 .
\]
Experiment agrees with $\tfrac{33}{9}$, \emph{not} $\tfrac{11}{9}$.
The factor of $3$ is direct, quantitative evidence that quarks carry
three colors.
}

\textit{Significance:} The number of colors is \emph{measured}, not
assumed. $R$ jumps by the right factor of $3$, and steps up at each new
quark threshold by $N_c Q_q^2$ --- a direct readout of both $N_c$ and
the quark charges.

\nt{Pitfall: when we ``observe quarks'' in the final state we never see
free quarks --- they are confined and immediately \emph{hadronize} into
color-neutral jets of hadrons. The cross section $\sigma(e^+e^-\to
\text{hadrons})$ is what is measured; the parton-level
$\sigma(e^+e^-\to q\bar q)$ is the theoretical object the $R$-ratio
compares to.}

%%==================================================================
\subsection{Classifying Weak Processes}\label{subsec:lec21:taxonomy}

\subsubsection*{Conceptual Background}

Weak interactions are conventionally divided by the kind of particles
participating in the $W/Z$ vertices:
\begin{itemize}
  \item \textbf{Leptonic:} the $W/Z$ couples \emph{only} to leptons.
  \item \textbf{Semileptonic:} both leptons \emph{and} quarks
        participate.
  \item \textbf{Non-leptonic:} \emph{only} quarks participate (no
        leptons).
\end{itemize}

\paragraph{Leptonic examples.}
Muon-type tau decay,
\begin{equation}
  \tau^- \to \nu_\tau\,\mu^-\,\bar\nu_\mu ,
  \label{eq:lec21:taudecay}
\end{equation}
where the $\tau^-$ emits a $W^-$ and becomes $\nu_\tau$, and the $W^-$
then produces $\mu^-\bar\nu_\mu$. Another is inverse muon decay,
$\nu_\mu\,e^- \to \mu^-\,\nu_e$. Only leptons appear at the vertices.

\paragraph{Semileptonic example: neutron beta decay.}
The neutron ($udd$) decays to a proton ($uud$) by converting one of its
down quarks: $d\to u\,W^-$, followed by $W^-\to e^-\bar\nu_e$:
\begin{equation}
  n\to p\,e^-\,\bar\nu_e ,
  \qquad \text{amplitude}\ \propto V_{ud}\ (\approx 1).
  \label{eq:lec21:betadecay}
\end{equation}
The two spectator quarks are unaffected; the quark-level transition
$d\to u$ carries the CKM element $V_{ud}$. A related crossing is
$\bar\nu_e\,p\to e^+\,n$, the reaction through which the electron
antineutrino was originally detected.

\paragraph{Non-leptonic examples.}
Here the $W$ connects only quarks. For the lambda hyperon
($\Lambda=uds$):
\begin{equation}
  \Lambda \to p\,\pi^- ,
  \qquad \text{amplitude}\ \propto V_{us}\ (\approx\lambda),
  \label{eq:lec21:lambdadecay}
\end{equation}
via $s\to u\,W^-$ with $W^-\to \bar u\,d$; the $uud$ recombine into a
proton and the $\bar u d$ into a $\pi^-$. Similarly the charged kaon
decay
\begin{equation}
  K^+ \to \pi^+\,\pi^0 ,
  \qquad \text{amplitude}\ \propto V_{us}.
  \label{eq:lec21:kaondecay}
\end{equation}

\textit{Significance:} The classification tells you immediately which
CKM elements and color factors enter a process, and hence its expected
rate. A strangeness-changing decay pays a factor $\lambda=|V_{us}|$ per
vertex; a same-generation transition pays nothing.

\nt{Pitfall (exam tip): you are \emph{not} expected to memorize the
quark content of exotic hadrons. You \emph{are} expected to draw the
quark-level diagram once the content is given, identify which quark
emits the $W$, and read off the CKM suppression of the active vertex.}

%%==================================================================
\subsection{Beyond the Standard Model}\label{subsec:lec21:bsm}

\subsubsection*{Physical Motivation}

The SM is extraordinarily successful, yet it is demonstrably incomplete.
This is where high-energy research begins. The professor listed the main
open questions:
\begin{itemize}
  \item \textbf{Neutrino masses.} The SM predicts $m_\nu=0$, but
        oscillations prove $m_\nu\neq0$. This is the one
        \emph{experimentally established} failure of the renormalizable
        SM.
  \item \textbf{Naturalness} (the hierarchy / fine-tuning problem) ---
        discussed below.
  \item \textbf{Dark matter.} Most of the matter in the universe is not
        made of any SM particle; its identity is unknown.
  \item \textbf{Baryogenesis.} Why is the universe made of matter and
        not antimatter? The SM cannot generate the observed asymmetry.
  \item \textbf{The flavor puzzle.} Why the observed pattern of fermion
        masses (huge hierarchies) and of CKM mixing? More a ``puzzle''
        than a sharp inconsistency.
\end{itemize}

\subsubsection*{The naturalness (hierarchy) problem}

The Higgs mass receives quantum corrections from loops --- for example a
top-quark loop. The physical (measured) mass is the sum of a bare term
and these corrections:
\begin{equation}
  \boxed{
    m_h^2\big|_{\rm phys}
    = m_{h,\rm bare}^2 + c\,\Lambda^2 ,
  }
  \label{eq:lec21:hierarchy}
\end{equation}
where $\Lambda$ is the cutoff (the scale up to which the SM is valid)
and $c=\mathcal O(1)$. The correction is \emph{quadratically divergent}
in $\Lambda$ --- the same quadratic divergence we met when discussing
loops earlier in the course.

If nothing new exists below the Planck scale,
$\Lambda\sim M_{\rm Pl}\sim 10^{18}~\text{GeV}$, then
$\Lambda^2\sim 10^{36}~\text{GeV}^2$, while the measured
$m_h^2\sim(125~\text{GeV})^2\sim 10^{4}~\text{GeV}^2$. The bare term must
then cancel the loop correction to $\sim$ $32$ decimal places --- an
extraordinary \emph{fine-tuning}. Equivalently: why is the weak force so
much stronger than gravity?

\textit{Significance:} A fundamental scalar at $125~\text{GeV}$, with its
mass sensitive to physics all the way up to $M_{\rm Pl}$, looks as
unnatural as a sharpened pencil balanced perfectly on its tip --- one
suspects a hidden mechanism (a ``string'') holding it up.

\nt{Pitfall: the hierarchy problem afflicts \emph{fundamental scalars}
only. Fermion masses are protected by an approximate \emph{chiral
symmetry} (the mass term vanishes in the symmetric limit, so corrections
are proportional to the mass itself and only logarithmically divergent).
This is why the electron mass is ``natural'' but the Higgs mass is not.}

\paragraph{Directions for a solution.}
Several broad strategies address naturalness:
\begin{itemize}
  \item \textbf{Compositeness:} the Higgs is not fundamental but a bound
        state of size $R$, so the effective cutoff is $\Lambda\sim 1/R$
        at a much lower scale.
  \item \textbf{Supersymmetry:} relate the scalar to a fermion via a
        symmetry; the fermion's chiral protection then ``leans over'' to
        protect the scalar mass, cancelling the quadratic divergence.
  \item \textbf{Cosmological relaxation:} the Higgs mass is not a static
        constant but a parameter that evolved dynamically over the
        history of the universe to its observed value.
\end{itemize}
The ultimate goal of the field is to find the \emph{true} Lagrangian of
nature --- one simpler and more complete than the SM, addressing all of
the above. A broad program of current and planned experiments is set to
test these ideas.

%%------------------------------------------------------------------
\subsection*{To Remember}

\begin{itemize}
  \item \textbf{Fermion masses:} $m_f = y_f\,v/\sqrt2$ from the Higgs
        vev; neutrinos massless in the SM (and this is wrong in nature).
  \item \textbf{Neutral current:}
        $\lagr_{\rm NC}=\frac{e}{s_Wc_W}\sum_i\bar\psi_i
        (T^3_i-\sin^2\theta_W Q_i)\slashed Z\psi_i$ ---
        universal, diagonal, chiral; no tree-level FCNC.
  \item \textbf{Charged current:} leptons flavor-diagonal; quarks carry
        $V_{\rm CKM}$, left-handed only, non-universal/non-diagonal.
        $\Gamma(W\to u_i\bar d_j)\propto 3|V_{ij}|^2$; top too heavy to
        appear.
  \item \textbf{Higgs:} coupling $\propto m_f$, diagonal; dominant decay
        $h\to b\bar b$ (top kinematically closed); no tree $hgg$ or
        $h\gamma\gamma$ (loop-induced, top-dominated).
  \item \textbf{Accidental symmetries:} $U(3)^5$ (Yukawas off) $\to$
        $U(1)_B\times U(1)_e\times U(1)_\mu\times U(1)_\tau$ (on);
        proton stable because it is the lightest $B$-charged state.
  \item \textbf{CKM:} hierarchical,
        $V\sim\left(\begin{smallmatrix}1&\lambda&\lambda^3\\
        \lambda&1&\lambda^2\\\lambda^3&\lambda^2&1\end{smallmatrix}\right)$,
        $\lambda\approx0.225$; $3$ angles $+$ $1$ CP phase.
  \item \textbf{Color count:}
        $R=\sigma(\to\text{had})/\sigma(\to\mu\mu)=N_c\sum_q Q_q^2$;
        below top, $R=3\cdot\frac{11}{9}=\frac{33}{9}$ confirms $N_c=3$.
  \item \textbf{Weak-process taxonomy:} leptonic / semileptonic /
        non-leptonic; read off CKM suppression per active quark vertex.
  \item \textbf{Beyond SM:} neutrino mass, naturalness
        ($m_h^2=m_{\rm bare}^2+c\Lambda^2$, $32$-digit tuning if
        $\Lambda\sim M_{\rm Pl}$), dark matter, baryogenesis, flavor.
\end{itemize}

%%------------------------------------------------------------------
\begin{furtherreading}
\subsection*{Further Reading}

\begin{itemize}
  \item D.\,J. Griffiths, \textit{Introduction to Elementary Particles}
        (2nd ed.): Ch.~9--10 for neutral/charged currents, the $R$-ratio,
        and the Cabibbo angle; \S2 for the weak-process classification.
  \item F. Halzen \& A.\,D. Martin, \textit{Quarks and Leptons}: Ch.~12--13
        for the electroweak currents and $e^+e^-\to$ hadrons.
  \item Y. Grossman \& Y. Nir, \textit{The Standard Model}: chapters on
        the SM currents, accidental symmetries, and baryon/lepton number
        --- the closest match to this course's logic and notation.
  \item M.\,E. Peskin \& D.\,V. Schroeder, \textit{An Introduction to QFT}:
        \S20 for the electroweak currents; \S5.3 for the $R$-ratio.
  \item For naturalness/BSM: G.\,F. Giudice, ``Naturally Speaking'';
        Particle Data Group, \textit{Review of Particle Physics} (2024),
        reviews on the Higgs boson and on physics beyond the SM.
        \url{https://pdg.lbl.gov/}
\end{itemize}
\end{furtherreading}
